% Generated draft: requires mathematical and source review.
\documentclass[11pt]{article}
\usepackage[margin=25mm]{geometry}
\usepackage{fontspec}
\setmainfont{Latin Modern Roman}
\usepackage{amsmath,amssymb,mathtools,unicode-math}
\setmathfont{Latin Modern Math}
\usepackage{enumitem,needspace,xurl,hyperref}
\setlength{\parindent}{0pt}
\setlength{\parskip}{5pt}
\setlength{\emergencystretch}{3em}
\newcommand{\N}{\mathbb{N}}
\newcommand{\Z}{\mathbb{Z}}
\newcommand{\Q}{\mathbb{Q}}
\newcommand{\R}{\mathbb{R}}
\newcommand{\C}{\mathbb{C}}
\newcommand{\F}{\mathbb{F}}
\newcommand{\A}{\mathbb{A}}
\newcommand{\PP}{\mathbb{P}}
\newcommand{\E}{\mathbb{E}}
\newcommand{\eps}{\varepsilon}
\newcommand{\dd}{\,\mathrm{d}}
\begin{document}
\section*{UnsolvedMath ID: 30000405}
\textbf{Source problem code:} catalogue entry
\subsection{Definitions and mathematical statement}
Let \(X,X'\) be smooth projective varieties over \(\Q\), and suppose their Hasse--Weil zeta functions, with the same local-factor conventions, satisfy
\[
\zeta_X^{\mathrm{HW}}(s)=\zeta_{X'}^{\mathrm{HW}}(s).
\]
At a prime of good reduction, the usual local factor is obtained from
\[
Z(X_p,T)=\exp\left(\sum_{r\geq1}\#X_p(\F_{p^r})\frac{T^r}{r}\right),
\]
with \(T=p^{-s}\); the treatment of bad primes and whether equality means a full or incomplete Euler product are not specified in the extract.

The source asks whether there is a common finite branched cover. In the literal algebraic reading, this means a variety \(Z\) with finite surjective morphisms
\[
Z\xrightarrow{\pi}X,\qquad Z\xrightarrow{\pi'}X',
\]
with ramification permitted. This is stronger than merely having two dominant rational maps from a resolution, but weaker than requiring both maps to be finite \emph{\'etale}. The source does not specify whether \(Z\) must be smooth, normal, geometrically integral, or defined over \(\Q\) rather than a field extension, nor any restrictions on branching or degrees. Those choices are essential to identify the intended theorem; they are not fixed by the phrase alone. A literal unrestricted branched-cover question must also be distinguished from a more restrictive commensurability question, and no stronger unramified condition is inserted silently.

On a specified common-cover convention, the proposed implication is therefore
\[
\forall X,X'\quad
\zeta_X^{\mathrm{HW}}=\zeta_{X'}^{\mathrm{HW}}
\ \Longrightarrow\ 
\exists(Z,\pi,\pi')\text{ satisfying that convention}.
\]
Connectedness and geometric-integrality conventions for the original varieties also require confirmation. A common finite cover by an integral \(Z\) would in particular have the same dimension as both targets; this is not an extra equality-of-dimension hypothesis supplied by the text.

The supplied literature assessment checked 2026-08-21 reports that the abelian-variety case follows from Faltings' isogeny theorem, while supplying no general theorem for all smooth projective varieties. An isogeny gives a specific finite-cover relationship in that reported case. The report's general partially solved label is not independently certified for every interpretation of unrestricted branched cover.
\subsection{Short English statement}
Does equality of Hasse--Weil zeta functions force a common finite cover with allowed branching? The source must specify the cover's field, smoothness, and ramification requirements; replacing branched by \emph{\'etale} would change the question.
\subsection{Sources}
\begin{itemize}
\item Local UnsolvedMath registry, permanent ID 30000405, catalogue entry, \emph{Common Branched Covers from Equal Hasse--Weil Zeta Functions}.
\item \emph{The Arithmetic of Fields} (2006), Oberwolfach Reports, pp. 317--378, DOI \url{https://doi.org/10.4171/owr/2006/06}.
\item The supplied assessment attributes the abelian-variety case to Faltings' isogeny theorem but gives no separate bibliographic citation or exact theorem statement.
\item Supplied literature triage checked 2026-08-21; it cites only the original report and supplies no precise general common-cover convention.
\end{itemize}
\end{document}
